% A.80: full root and independent proof/source/code audits PASS.
\subsection{A genus-two intermediate with a noncyclic derived closure}

\begin{statement}{A.80}{Proved}{The natural genus-two wreath cover has the required fixed groups}
Choose \(s\) with \(s^4+2s^2+2=0\), and put
\[
 i=-1-s^2,\qquad q=3+2s+2s^2,\qquad
 E_4:y^2=x^3-6x^2+x,\qquad Q=(q,-2sq).
                                                        \tag{A.80.1}
\]
Then \(2Q=R=(1,2i)\), and \(R\) has order four.
Let \(Y\) be the smooth proper normalization of
\[
 v^2=(u^2+q)^3-6(u^2+q)^2+(u^2+q).
                                                        \tag{A.80.2}
\]
Its actual map to \(E_4\) is
\(q_Y(u,v)=(u^2+q,v)+Q\). Define
\[
 \phi(x,y)=\left(\frac{(x-1)^2}{4x},
                  \frac{y(1-x^2)}{8x^2}\right),\qquad
 \eta=[1+i]\circ\phi:E_4\longrightarrow E,
 \quad E:y^2=x^3-x.
                                                        \tag{A.80.3}
\]
The map \(f=\eta q_Y\) has degree eight and is ramified only
above \(O\). Its geometric Galois closure has group
\(C_2\wr C_4=(C_2)^4\rtimes C_4\), of order sixty-four.
The affine curve \(X=Y\setminus f^{-1}(O)\) has genus two
and six cusps, with ramification indices \(2,2,1,1,1,1\).

For every original packet of Part A, every coefficient prime
\(\ell\), and every permitted open \(\Delta_0\) preserving
this actual cover and the selected cusp lift \(P\),
\[
       \pi_1(X_{\overline K})^{(\ell),s_b(\Delta_0)}
                              =I_{P,\ell}.              \tag{A.80.4}
\]
All its pro-\(\ell\) normalizer assertions follow from A.59.
At odd \(\ell\), the fixed-group proof also holds for every
packet at any cusp over a finite number field of definition,
with that field's closed derived subgroup or any open subgroup
of it. This statement concerns this actual genus-two intermediate;
it does not assert fixedness for every intermediate of the
genus-seventeen Galois closure.
\end{statement}

\begin{proof}
\paragraph{The actual elliptic isogeny and wreath closure.}
The displayed identities give \(i^2=-1\), \(2Q=R\),
and \(2R=(0,0)=T\). The rational functions in (A.80.3)
satisfy the equation of \(E\), define a degree-two isogeny
with kernel \(\{O,T\}\), and send \(R\) to \((0,0)\).
Consequently \(\eta\) is a cyclic degree-four isogeny with
\[
                     \ker\eta=K=\{O,R,T,-R\}.
                                                        \tag{A.80.5}
\]
These are identities of actual maps, not only identities of
elliptic \(j\)-invariants.

Let \(\ell_Q\) be the tangent line to \(E_4\) at \(Q\).
The function \(a=(x-1)/\ell_Q\) has divisor
\(R+O-2Q\). The connected double cover \(z^2=a\) is therefore
branched exactly at \(O,R\) and has genus two. On translating
the elliptic variable by \(-Q\), the function \(x-q\)
has the same divisor as \(a\). Their ratio is a nonzero
constant, which can be made a square after a finite number-field
extension. This identifies its smooth normalization with
(A.80.2) and its actual elliptic map with \(q_Y\).
The two remaining points of \(K\) each have two unramified
preimages, proving the stated six-cusp description.

For the full Galois closure, take the four translates
\(a_j(P)=a(P+jR)\), \(0\le j<4\). Their square classes
are linearly independent. Indeed the kernel of the cyclic
map \(\mathbb F_2[C_4]\to\langle[a_j]\rangle\) is an
ideal of \(\mathbb F_2[t]/((t+1)^4)\). Every nonzero ideal
contains \((t+1)^3=1+t+t^2+t^3\). A nonzero kernel would
make \(\prod_j a_j\) a square. But half its divisor is
\(D_K-D_{Q+K}\), whose degree-zero class on \(E_4\) is
\(-4Q=-2R=T\ne0\). Thus that product is not a square.
Adjoining the four roots gives the full \((C_2)^4\)-extension;
translation by \(R\) lifts by cyclically permuting the four
roots, literally with order four. The group over \(E\) is
therefore the stated wreath product. Its natural degree-eight
action is on four pairs; the cusp commutator exchanges one
pair in each of two adjacent blocks. This gives the displayed
cycle lengths. Riemann--Hurwitz gives genus seventeen and
thirty-two cusps for the full closure.

\paragraph{The actual complementary elliptic factor is non-CM.}
Put
\[
 A=3q-6,\qquad B=3q^2-12q+1,\qquad C=q^3-6q^2+q.
                                                        \tag{A.80.6}
\]
The two involutions \(\sigma(u,v)=(-u,v)\) and
\(h(u,v)=(u,-v)\) give the actual elliptic quotients
\(E_4\) and
\[
 F:\quad V^2=t(t^3+At^2+Bt+C),\qquad t=u^2,\quad V=uv.
                                                        \tag{A.80.7}
\]
The second quotient is by \(\sigma h\). The change
\(X=C/t,\ Y_F=CV/t^2\) gives
\[
             Y_F^2=X^3+BX^2+ACX+C^2.                  \tag{A.80.8}
\]
The two quotient pullbacks are the \(\sigma\)-eigenspaces
of proper homology. They have dimension two each, and give
an actual isogeny \(J_Y\sim E_4\times F\). In particular
the complementary factor and its projector are defined by
the displayed maps.

Here is an exact ordinary-reduction certificate for non-CM.
At the degree-one primes of \(\mathbb Q(s)\) specified below,
(A.80.8) has the following reductions:
\[
\begin{array}{c|c|c|c|c|c}
 p&s\bmod p&(B,AC,C^2)\bmod p&\#F(\mathbb F_p)&a_p&a_p^2-4p\\\hline
29&4&(15,25,24)&34&-4&-100\\
37&17&(17,21,33)&42&-4&-132
\end{array}                                                   \tag{A.80.9}
\]
Both discriminants of the cubic models are nonzero, and both
reductions are ordinary since \(p\nmid a_p\). The exact
Legendre-symbol sums for every field element are retained in
the certificate. By \cite[Theorem V.3.1, printed p.145]
{fghSilvermanAEC}, their geometric rational endomorphism
algebras are respectively \(\mathbb Q(i)\) and
\(\mathbb Q(\sqrt{-33})\). The Frobenius eigenvalue ratio
has nonzero valuation at a prime above \(p\), because the
ordinary slopes are zero and one; it is not a root of unity.
Thus every positive Frobenius power still generates the same
quadratic field. If \(F\) had CM, specialization after a
finite extension defining its endomorphisms would embed the
same quadratic CM field into both of these distinct quadratic
fields, which is impossible. Hence \(F\) is non-CM.
This proof does not infer non-CM from a numerical approximation
to \(j\), nor from a conjectural classification of CM points.

\paragraph{The proper compact-type model and its actual map.}
Work at residue characteristic two. The wreath cover is an
actual one-point Galois two-group cover of the potential-good
supersingular elliptic curve \(E\). A.78 gives potential
good Jacobian reduction with special two-rank zero for its
full closure. The norm onto \(J_Y\) is surjective; hence
the same properties hold for \(J_Y\). After finite extension,
the proper stable model of \(Y\) has a tree, and every
positive-genus component has two-rank zero. Its only possible
proper stable shapes are one smooth genus-two component or
two supersingular elliptic components joined at a node.
Indeed a stable rational vertex has valence at least three,
whereas a tree with total positive genus two has at most two
positive-genus leaves.

The actual map \(q_Y\) extends to the good elliptic model
of \(E_4\), by the extended Abel--Jacobi map and the
homomorphism of good Jacobians; see
\cite[Theorem 1.1 and Proposition 3.2]{fghDupuyRabinoff}.
The degree of the pullback of a degree-one line bundle is
two on each fibre. Rational components map constantly, so
the sum of the degrees on positive-genus components is two.
All points of \(K\) specialize to \(O\), since the good
elliptic special fibre is supersingular and \(K\) is
two-primary. The automorphism \(\sigma\) extends to the
proper stable model. It cannot act identically on every
positive-genus component: the good-Jacobian Tate comparison
would then make it the identity on proper homology, whereas
the actual quotient decomposition above has a nonzero
\(-1\)-eigenspace \(V_F\). This supplies the precise
specialization faithfulness needed below.

If the special fibre is smooth of genus two, its degree-two
map to the elliptic target is separable. A purely inseparable
degree-two map between smooth curves over the algebraically
closed residue field would preserve genus. The two generic
ramification cusps specialize to fixed points of \(\sigma\).
Deuring--Shafarevich, in
\cite[Corollary 1.8, printed pp.35--36]{fghCrew1984},
with source and quotient of two-rank
zero, gives exactly one such point. Its image is \(O\),
and it is the unique point of the special fibre above \(O\).
Thus all six cusps specialize there. After separating the
marks, their connecting paths contain only rational components.

If the two component degrees are \((1,1)\), \(\sigma\)
exchanges the components. Otherwise its restriction to each
would be the identity, since each degree-one map is an
isomorphism, contradicting faithful specialization. The
two ramification cusps specialize to the node, so its image
is \(O\). Each elliptic component has just that point
above \(O\). Again all marks lie in the rational tree
inserted over that node. In both cases A.41 applies at every
odd coefficient prime and to every selected cusp.

The remaining case has degrees \((2,0)\). Denote the main
component by \(B_E\) and the contracted tail by \(B_F\).
The map \(B_E\to(E_4)_s\) is radicial: a degree-two
separable isogeny from a supersingular elliptic curve does
not exist. Therefore \(\sigma\) is the identity on
\(B_E\), and acts nontrivially on \(B_F\). It fixes
their node, and this is its unique fixed point on the
supersingular elliptic tail, again by Deuring--Shafarevich.

If the image of the node is different from \(O\), all six
cusps specialize to the single point of \(B_E\) over
\(O\), and A.41 applies. Suppose instead that the node
maps to \(O\). The two ramification cusps must specialize
to that node, since they are fixed by \(\sigma\). Any
selected cusp with the same proper specialization has a
rational companion path. A remaining selected cusp can
specialize smoothly on \(B_F\). Choose one of the two
ramification cusps as \(Q_0\). In the pointed stable
tree the path from the selected cusp to \(Q_0\) avoids
the positive-genus vertex \(B_E\). The vertex carrying
\(Q_0\) lies in the rational tree inserted over the old
node, not on either elliptic component. Thus the distance
from \(B_E\) to that path is strictly positive.

For clarity, the local handle blocks in this last configuration
are the actual global elliptic factors. The specialization of
\((q_Y)_*\) is a nonzero isogeny on \(B_E\) and zero on
\(B_F\). Its adjoint pullback therefore identifies the
global \(V_E=V_\ell E_4\) with the first local block and
the global Prym \(V_F=V_\ell F\) with the second. This
uses the actual map and its polarization adjoint. It does
not apply arbitrary global endomorphisms to a local fibre.

\paragraph{A partial-distance kernel in every weight.}
We prove the precise extension of the calculation in A.41
needed here. Let a compact-type pointed graph have two
handle blocks \(V_E,V_F\). Root it at the selected cusp
\(P\), eliminate the generator of a distinct cusp \(Q_0\),
and retain the other cusp generators \(T\), including
\(t_P\). Give handles weight one and cusps weight two.
Assume only that \(B_E\) is off the \(P\)--\(Q_0\)
path. No condition is imposed on the location of \(B_F\).
The actual first weighted graph derivation, including the
arbitrary packet term, has the form
\[
 Dv=[W_E,v]\ (v\in V_E),\quad
 Dv=[W_F,v]\ (v\in V_F),\quad
 Dt_j=[W_j,t_j],
                                                        \tag{A.80.10}
\]
where each \(W\) is a linear combination of
\(\Omega_E,\Omega_F,T\). Its coefficient of \(\Omega_E\)
in \(W_E\) is
\[
 r_E=\sum_e m_e E_e(1-Q_{0,e})>0.                       \tag{A.80.11}
\]
This is exactly the positive distance just established.
Adding \(\lambda t_P\) to every \(W\), for arbitrary
\(\lambda\in\mathbb Q_\ell\), does not change it.
The based comparison and the absence of a weight-one inner
discrepancy are precisely A.41.4--A.41.6, using
\cite[Definition 3.1.5 and Theorem 3.1.6]{fghBettsDogra}.

Put \(k=\mathbb Q_\ell\), and inside the free associative
enveloping algebra set
\[
             B=k\langle T,V_F,\Omega_E\rangle.
                                                        \tag{A.80.12}
\]
These are free primitive generators. The component algebra
\(k\langle V_E\rangle\) has the triangular bimodule
basis A.41.8 over \(k[\Omega_E]\); take its free product
with the free algebra on \(T,V_F\). This gives a normal
basis in which every word outside \(B\) has a first
nonboundary \(E\)-handle core preceded by a \(B\)-word.
Unlike A.41, the individual \(F\)-handle letters belong
to the coefficient algebra; this change is essential.

The derivation preserves \(B\): \(\Omega_F\) already
belongs to the algebra generated by \(V_F\), and
\(D\Omega_E=[W_E,\Omega_E]\). For a fixed first core
and its following tail, the operator which raises the
weight of its initial coefficient slot by two is
\[
                         K=D_B+R_{W_E}.                \tag{A.80.13}
\]
All terms acting after that core preserve the initial-slot
weight. They cannot remove the core; the triangular
bimodule basis retains end powers of \(\Omega_E\),
and any other coefficient letter separates components.
Thus the top initial-slot operator acts independently on
each fixed core tail. We use weighted slot degree here,
since \(F\)-handle letters have weight one.

Count occurrences of the single free letter \(\Omega_E\)
in \(B\). The part of \(K\) increasing that count is
\[
       K^+=D_c+r_E R_{\Omega_E},\qquad
       D_c\Omega_E=0,\qquad
       D_c z=c_z[\Omega_E,z]\quad(z\in T\cup V_F).
                                                        \tag{A.80.14}
\]
All omitted terms preserve this count. Fix the word
\(z_1\cdots z_n\) in the other free letters and insert
powers of \(\Omega_E\) in its \(n+1\) slots. The
resulting polynomial module \(k[Z_0,\ldots,Z_n]\) has
\(K^+\) equal to multiplication by
\[
 c_{z_1}Z_0+
 \sum_{h=1}^{n-1}(c_{z_{h+1}}-c_{z_h})Z_h
                  +(r_E-c_{z_n})Z_n.                  \tag{A.80.15}
\]
For \(n=0\) this is \(r_EZ_0\). In every case the
coefficients sum to \(r_E\ne0\), so the polynomial is
nonzero and multiplication is injective. The maximum
\(\Omega_E\)-count of a finite input proves that \(K\)
is injective. The maximum initial-slot weight then proves,
in every weighted degree,
\[
       \ker D\subset B\cap\mathbb L(V_E\oplus V_F\oplus T)
                       =\mathbb L(T,V_F,\Omega_E).
                                                        \tag{A.80.16}
\]
The last equality takes primitive elements in the free
tensor Hopf subalgebra (A.80.12). This proves an unbounded
kernel inclusion; no finite-degree rank extrapolation is used.

\paragraph{The non-CM block makes every surviving fixed grade Tate.}
Choose a finite number field \(L\) defining all maps,
cusps and the model, with connected algebraic monodromy
and \(\Gamma=\overline{[G_L,G_L]}\) contained in the
prescribed open derived subgroup. Put
\(H=\pi_1(X_{\bar L})^{(\ell),s(\Gamma)}\), and take
the Lie algebra \(\mathfrak h\) of the unipotent closures
of this actual subgroup, with its induced weight filtration.
It is stable under \(G_L\), whose action on it is abelian.
The actual local Frobenius--tame-inertia commutator in
\(\Gamma\) has first weighted derivation a nonzero common
multiple of (A.80.10). Hence (A.80.16) bounds the graded
pieces of \(\mathfrak h\), exactly as in A.41.

The global proper representation is the actual direct sum
\(V_E\oplus V_F\). The first factor has CM over a finite
extension; the second is non-CM by (A.80.9). Serre's open
image theorem, \cite[Theorem 3, printed p.299]
{A_Serre1972}, shows that the connected derived algebraic
monodromy group is
\[
                         \{1\}\times\mathrm{SL}(V_F).
                                                        \tag{A.80.17}
\]
Indeed its projection to the CM torus is trivial and its
projection to \(\mathrm{GL}(V_F)\) has derived group
\(\mathrm{SL}(V_F)\). The actual \(\Gamma\)-fixed
graded pieces are therefore invariant under that full
\(\mathrm{SL}_2\).

In (A.80.16), every generator of \(T\) and \(\Omega_E\)
has weight two and character \(\chi_\ell\). A tensor
summand with \(r\) letters in \(V_F\) and \(a\) of
these Tate generators has weight \(r+2a\). Its
\(\mathrm{SL}(V_F)\)-invariants vanish for odd \(r\);
for even \(r\), \(\mathrm{GL}(V_F)\) acts on them
by the scalar \(\det^{r/2}\). One may see this directly:
the action factors through \(\mathrm{GL}_2/\mathrm{SL}_2\),
and a scalar matrix acts on the tensor power by its
\(r\)-th power. Since \(\det V_F=\chi_\ell\), the
entire invariant summand has scalar character
\(\chi_\ell^{a+r/2}\). The same holds for its primitive
Lie subspace. Consequently
\[
 \operatorname{gr}_{2m+1}\mathfrak h=0,\qquad
 G_L\text{ acts on }\operatorname{gr}_{2m}\mathfrak h
                           \text{ by }\chi_\ell^m.
                                                        \tag{A.80.18}
\]
Apply the actual-subgroup cyclotomic criterion of A.38.
It gives \(H=I_{P,\ell}\), with no lifting of arbitrary
rational Lie elements. This handles the only configuration
not already covered by rational companion paths and proves
the odd-prime assertions for every cusp packet.

For the original Part A packets at \(\ell=2\), the actual
wreath closure is a two-group, so A.22 gives the required
open embedding and exact fixed inertia. Normalizers at all
primes follow from A.59. Arithmetic comparisons use actual
cover maps after finite field extension and geometric paths
only up to inner conjugacy; no arbitrary profinite path is
asserted to be defined over a finite extension.

\paragraph{Exact certificates and scope.}
The retained programs
\path{code/a_wreath64_frontier.py},
\path{code/a_wreath64_map.py}, and
\path{code/a_wreath64_prym.py}
check the finite group/action, rational map identities,
quotient equations and both complete finite-field counts.
Their JSON and stdout are retained; the read-only aggregate replay is
\path{code/a_wreath64_genus_two.py}. The additional exact
\(j\)-polynomial and Sturm sequence are corroboration only;
the non-CM proof uses the two ordinary fields above.
The graph and all-weight kernel arguments are proved here,
independently of finite algebraic checks. The full closure
and its other intermediate fixed groups remain outside
this statement.
\end{proof}
